\documentclass[12pt]{article}
\usepackage[letterpaper,margin=1in]{geometry}
\usepackage{amsmath,amssymb,setspace}
\usepackage{graphicx}
\usepackage{bm}
\usepackage{tcolorbox}
\usepackage{enumitem}
\usepackage{xcolor}
\usepackage{hyperref}
\hypersetup{
    colorlinks=true,
    linkcolor=blue,
    urlcolor=cyan,
}
\begin{document}
\doublespacing

\title{Solutions to Homework 03 -- Quiz Questions}
\author{}
\date{}
\maketitle

\section*{Question 1}

The problem provides a labeled data set
\[
\mathcal{D} = \{(\bm{x}_i, c_i) \mid i = 1, 2, \ldots, N\},
\]
where each \(\bm{x}_i \in \mathbb{R}^D\) is a feature vector and \(c_i \in \{0,1\}\) is the class label (with \(0\) for no damage and \(1\) for damaged). The model is expressed as
\[
m(\bm{x}\mid \bm{w}),
\]
with parameters \(\bm{w}\).

\textbf{Answer:} \textbf{A. Discriminative Machine Learning}

\textbf{Explanation:} The model directly learns a mapping from \(\bm{x}\) to a score (and then a decision) without modeling the full joint distribution \(p(\bm{x}, c)\). This is typical of discriminative approaches (e.g., logistic regression, SVMs) where one estimates \(p(c \mid \bm{x})\) or the decision boundary directly.

\bigskip

\section*{Question 2}

The decision function is given by
\[
\text{if } f(\bm{x}\mid \bm{w}) \geq 1 \text{ then } c = 1,\quad \text{otherwise, } c = 0.
\]
Geometrically, the function’s contour \(f(\bm{x}\mid \bm{w}) = 1\) defines a boundary. The text explains that if a data point lies inside the contour, the structure is intact (i.e., no damage), and if outside, it is damaged.

For the testing data with computed values \((0.1,\; 1.5,\; 0.2,\; 2)\):
- A score of \(0.1\) is \textbf{inside} the contour (\(0.1 < 1\)) \(\rightarrow\) \textbf{no damage} (\(c=0\)).
- A score of \(1.5\) is \textbf{outside} the contour (\(1.5 \ge 1\)) \(\rightarrow\) \textbf{damaged} (\(c=1\)).
- A score of \(0.2\) is \textbf{inside} the contour (\(0.2 < 1\)) \(\rightarrow\) \textbf{no damage} (\(c=0\)).
- A score of \(2\) is \textbf{outside} the contour (\(2 \ge 1\)) \(\rightarrow\) \textbf{damaged} (\(c=1\)).

\textbf{Answer:} The predicted states are:
\[
\text{no damage},\quad \text{damaged},\quad \text{no damage},\quad \text{damaged}.
\]

\bigskip

\section*{Question 3}

There are two parts here:

\begin{enumerate}[label=(\alph*)]
    \item If the decision function is defined by a fixed number of parameters, here \(\bm{w}\) is a \(10\times 1\) vector, then the model has a fixed (finite) number of parameters.
    
    \textbf{Answer (a):} \underline{parametric} model.
    
    \item If instead the engineer develops a model where the number of parameters scales with the number of training samples (e.g., using a kernel-based method), then the model adapts its complexity to the data size.
    
    \textbf{Answer (b):} \underline{non-parametric} model.
\end{enumerate}

\bigskip

\section*{Question 4}

The engineer splits the data into
\[
D_0 = \{(x_i, c_i = 0) \mid i = 1, 2, \ldots, N_0\} \quad \text{and} \quad D_1 = \{(x_i, c_i = 1) \mid i = 1, 2, \ldots, N_1\},
\]
and then uses a Gaussian mixture model to model \(p(x\mid c=0)\) and \(p(x\mid c=1)\) separately. The priors \(P(c=0)\) and \(P(c=1)\) are also estimated from the data.

\textbf{Answer:} \textbf{A. Generative machine learning}

\textbf{Explanation:} In this approach the engineer is modeling the class-conditional densities \(p(x\mid c)\) and the priors, which is the hallmark of generative methods.

\bigskip

\section*{Question 5}

We are given:
\[
P(c=0)=0.99,\quad P(c=1)=0.01.
\]
And for the two feature vectors, the likelihoods are:
\[
P(x_1\mid c=0)=0.5,\quad P(x_1\mid c=1)=0.99,
\]
\[
P(x_2\mid c=0)=0.001,\quad P(x_2\mid c=1)=0.99.
\]

Using Bayes' theorem, we compute the unnormalized posterior probabilities.

\textbf{For \(\bm{x}_1\):}
\[
P(c=0 \mid x_1) \propto P(x_1 \mid c=0)P(c=0)=0.5 \times 0.99=0.495,
\]
\[
P(c=1 \mid x_1) \propto P(x_1 \mid c=1)P(c=1)=0.99 \times 0.01=0.0099.
\]
Since \(0.495 \gg 0.0099\), the posterior probability favors \(c=0\).

\textbf{Answer:} \(\bm{x}_1\) belongs to the \underline{no damage} class.

\bigskip

\textbf{For \(\bm{x}_2\):}
\[
P(c=0 \mid x_2) \propto P(x_2 \mid c=0)P(c=0)=0.001 \times 0.99=0.00099,
\]
\[
P(c=1 \mid x_2) \propto P(x_2 \mid c=1)P(c=1)=0.99 \times 0.01=0.0099.
\]
Here, \(0.0099 \gg 0.00099\), so the posterior probability favors \(c=1\).

\textbf{Answer:} \(\bm{x}_2\) belongs to the \underline{damaged} class.

\bigskip

\section*{Summary of Answers}

\begin{itemize}
    \item \textbf{Q1:} A. Discriminative Machine Learning.
    \item \textbf{Q2:} The predicted states for scores \((0.1,\,1.5,\,0.2,\,2)\) are: no damage, damaged, no damage, damaged.
    \item \textbf{Q3:} (a) \underline{parametric} model; (b) \underline{non-parametric} model.
    \item \textbf{Q4:} A. Generative machine learning.
    \item \textbf{Q5:} \(\bm{x}_1\) is \underline{no damage} and \(\bm{x}_2\) is \underline{damaged}.
\end{itemize}

\end{document}
